\documentclass[10pt]{article}

\usepackage[utf8]{inputenc}

\usepackage[T1]{fontenc}

\usepackage{amsmath}

\usepackage{amsfonts}

\usepackage{amssymb}

\usepackage[version=4]{mhchem}

\usepackage{stmaryrd}

\usepackage{graphicx}

\usepackage[export]{adjustbox}

\graphicspath{ {./images/} }

\usepackage{bbold}



\begin{document}

$\Omega-\omega$, к-рые соответственно называются критической и допустимой. Считают, что при попадании проверочной статистики $X$ в критич. область $\omega$ гипотеза $H_{0}$ неверна (верна $H_{1}$ ), а при попадании $X$ в допустимую область гипотеза $H_{0}$ верна ( $H_{1}$ ошибочна).



Разделение пространства $\Omega$ на критическую и допустимую области обычно производится так, чтобы вероятность отвергнуть гипотезу, когда она верна (то есть вероятность потери), была бы малой. Величину этой вероятности называют у р о в не м знач и м ос т и, или величиной критерия. Таким образом, уровень значимости $\alpha$ равен вероятности попадания $X$ в $\omega$, когда гипотеза $H_{0}$ верна, то есть



\[

\mathrm{P}\left\{X \in \omega \mid H_{0}\right\}=\alpha .

\]



С другой стороны, целесообразно потребовать также малости вероятности принятия ложной гипотезы, то есть в е роятности примеси $\beta$ :



\[

P\left\{X \in \Omega-\omega \mid H_{1}\right\}=\beta

\]



Для оценки критерия проверки альтернативных гипотез служит величина, называемая мош нос т ю к р и те р и я, к-рая определяется как вероятность $1-\beta$ попадания $X$ в критич. область пространства $\Omega$, когда верна гипотеза $H_{1}$, то есть



\[

\mathrm{P}\left\{X \in \omega \mid H_{1}\right\}=1-\beta .

\]



При выборе гипотезы исследователь обычно решает, какие потери $\alpha$ он может допустить, а затем выбирает проверочную статистику и критич. область так, чтобы максимизировать мощность критерия $1-\beta$.



Одна из наиболее общих проверяемых гипотез при А. д. состоит в том, что плотность вероятности $p(\boldsymbol{x})$ есть данная функция $\boldsymbol{x}$, то есть $p(\boldsymbol{x})=f(\boldsymbol{x})$. Здесь обычно нет определенной альтернативной гипотезы, то есть фактически имеется набор всевозможных альтернативных гипотез, к-рые явно не определены. В этом случае невозможно вычислить примесь и определить мощность критерия. Такая задача возникает при проверке совпадения экспериментальных данных с какой-либо теоретич. моделью и решается на основе кри т е рия с оглас ия. Как при обычной проверке гипотез, начинают с выбора проверочной статистики, однако пространство $\Omega$ не делится на критич. и допустимую области. Уровень значимости здесь определяется как вероятность того, что при условии $H_{0}$ проверочная статистика $X$ будет иметь значение, превышающее величину $T$, наблюдаемую из данных,



\[

\mathrm{P}\left\{X \geqslant T \mid H_{0}\right\}=\alpha(T) \text {. }

\]



В данном контексте величина $\alpha(T)$ называется также уровнем достоверности.



Критерий согласия конструируется при помощи меры различия между непараметрич. оценкой плотности вероятности и теоретич. функцией плотности вероятности проверяемой гипотезы. Наиболее популярной является квадратич. мера, нормированная на дисперсию. В достаточно общих предположениях проверочная статистика сводится к сумме квадратов независимых, нормально распределенных случайных величин с нулевым средним и единичной дисперсией, к-рая имеет $\chi^{2}$-распределение с числом степеней свободы, равным количеству членов в сумме. В этом случае критерием согласия является $\chi^{2}-$ критерий Пирсона.



Современные экспериментальные исследования в области ядерной физики, геофизики, физики атмосферы, океана и др. характеризуются огромным объемом получаемой первичной информации (до $10^{12}$ бит/с и более). Результаты эксперимента обычно составляют $\sim 10^{3}$ бит/с. Таким образом, в процессе А. д. происходит значительное сжатие информации (в 1 млрд. раз и более). А. Д. таких экспериментальных исследований немыслим без использования средств автоматизации и быстродействующей вычислительной техникИ.



Jum.: [1] М и т ро по льск и й А. К., Техника статистических вычислений, 2 изд., М., 1971; [2] Статистические методы в экспериментальной физике, пер. с англ., М., 1976; [3] Т ь юк и Дж., Ана- лиз результатов наблюдений, пер. с англ., М., 1981; [4] М о с те лле р Ф., Т ь юк и Дж., Анализ данных и регрессия, пер. с англ.,



\begin{center}

\includegraphics[max width=\textwidth]{2023_07_08_62a161bf9cc1ff3ff4c1g-1}

\end{center}



АНАЛИТИЧЕСКАЯ ПЕРЕНОРМИРОВКА - см. Перенормировка в квантовой механике.



АНАЛИтИЧЕСКАЯ РЕГУЛЯЦИЯ - см. Расходимость в квантовой теории поля, Ультрафиолетовая расходимость.



АНАЛИТИЧЕСКАЯ ФУНКЦИЯ, го л о о о ф н а я Фун к ц и я,- функция, которую можно разложить в степенной ряд. Класс А. Ф. охватывает большинство функций, встречающихся в основных вопросах математики и физики; основные операции арифметики, алгебры и анализа не выводят за пределы класса А. Ф.; наконец, А. Ф. обладают свойством единственности, то есть каждая А. Ф. образует одно органически связное единое целое во всей своей естественной области определения.



Пусть $D$ - область в плоскости комплексного переменного $\mathbb{C}$ (то есть открытое связное множество точек $\mathbb{C}$ ). Если каждой точке $z=x+i y \in D$ соответствует комплексное число $\boldsymbol{w}=\boldsymbol{u}+i v$, то в области $D$ определена (од нознач ная) функция $w=f(z)$ (или $f: D \rightarrow \mathbb{C}$ ) комплексного переменного $z$. Иначе говоря, функцию $w=f(z)=f(x+i y)$ можно рассматривать как комплексную функцию двух действительных переменных $x$ и $y$, определенную в области $D \subset \mathbb{R}^{2}$ ( $\mathbb{R}^{2}-$ евклидова плоскость). Или, еще, задание такой комплексной функции $f$ равносильно заданию двух действительных функций



\[

u=\varphi(x, y), v=\psi(x, y),(x, y) \in D .

\]



При фиксированном $z \in D$ каждому прирашению $\Delta z=\Delta x+i \Delta y, z+\Delta z \in D$, соответствует прирашение функции



\[

\Delta f(z)=f(z+\Delta z)-f(z)=\Delta u(x, y)+i \Delta v(x, y) .

\]



Если при $\Delta z \rightarrow 0$



\[

\Delta f(z)=A \Delta z+o(\Delta z)

\]



или, что то же, если существует предел



\[

\lim _{\Delta z \rightarrow 0} \frac{\Delta f(z)}{\Delta z}=A

\]



то функция $f$ называется д и ф ф е р е н ц и р у е м о й в смысле $\mathbb{C}$ (или в смысле комплексного анализа) в точке $z$ и $A=f^{\prime}(z)-$ п р о и з во д н а я функции $f$ в точке $z$, a $A \Delta z=f^{\prime}(z) d z(d z=\Delta z)-$ дифф е р е нци ал $f$ в точке $z$. Функция $f$ д иффе р е нци р у е м в области $D$, если она дифференцируема в каждой точке $D$.



Если функция $f$ рассматривается как функция двух действительных переменных (то есть в смысле $\mathbb{R}^{2}$ ), то ее дифференциал имеет вид



\[

d f=\frac{\partial f}{\partial x} d x+\frac{\partial f}{\partial y} d y

\]



где



\[

\frac{\partial f}{\partial x}=\frac{\partial \varphi}{\partial x}+i \frac{\partial \psi}{\partial x}, \quad \frac{\partial f}{\partial y}=\frac{\partial \varphi}{\partial y}+i \frac{\partial \psi}{\partial y}

\]



\begin{itemize}

  \item частные производные функции $f$ в точке $z$. Здесь можно перейти от независимых переменных $x, y$ к переменным $z=x+i y, \bar{z}=x-i y$, то есть записать $f$ в виде $f(z, \bar{z})$, причем тогда

\end{itemize}



\[

\begin{array}{cc}

x=\frac{z+\bar{z}}{2}, & y=\frac{z-\bar{z}}{2}, \\

d x=\frac{d z+d \bar{z}}{2}, & d y=\frac{d z-d \bar{z}}{2},

\end{array}

\]



и получается запись дифференциала $d f$ в комплексной форме



\[

d f=\frac{\partial f}{\partial z} d z+\frac{\partial f}{\partial \bar{z}} d \bar{z}

\]



где



\[

\frac{\partial f}{\partial z}=\frac{1}{2}\left(\frac{\partial f}{\partial x}-i \frac{\partial f}{\partial y}\right), \quad \frac{\partial f}{\partial \bar{z}}=\frac{1}{2}\left(\frac{\partial f}{\partial x}+i \frac{\partial f}{\partial y}\right)

\]



\begin{itemize}

  \item (формальные) производные $f$ по $z$ и $\bar{z}$.

\end{itemize}



АНАЛИЗ ДАННЫХ 21





\end{document}